% Part: first-order-logic
% Chapter: syntax-and-semantics
% Section: Structures

\documentclass[../../../include/open-logic-section]{subfiles}

\begin{document}

\olfileid{fol}{syn}{str}

\olsection{一阶语言的结构}

\begin{explain}
一阶语言本身是\emph{未解释的}：常项符号、函数符号和谓词符号没有特定的意义。
意义是通过指定一个\article{structure} \emph{结构}给出的。
该结构规定了\emph{论域}，即常项符号所指称的对象、函数符号所作用的对象，以及量词所遍历的对象。
此外，它还具体规定了哪个常项符号指称哪个对象、函数符号如何将对象映射到对象，以及谓词符号适用于哪些对象。
结构是逻辑中\emph{语义}概念（例如后承、有效性、可满足性）的基础。
它们在文献中有不同的称谓，如“结构”“解释”或“模型”。
\end{explain}

\begin{defn}[结构]
\Article{structure} 一阶逻辑语言 $\Lang{L}$ 的一个\emph{结构}~$\Struct M$ 由以下要素组成：
\begin{enumerate}
\item \emph{论域：}一个非空集合 $\Domain M$。
\item \emph{常项符号的解释：}对 $\Lang{L}$ 中的每个常项符号~$c$，指定一个元素 $\Assign{c}{M} \in \Domain M$。
\item \emph{谓词符号的解释：}对 $\Lang{L}$ 中的每个 $n$ 元谓词符号~$R$（$\eq$ 除外），指定一个 $n$ 元关系 $\Assign{R}{M} \subseteq \Domain{M}^n$。
\item \emph{函数符号的解释：}对 $\Lang{L}$ 中的每个 $n$ 元函数符号~$f$，指定一个 $n$ 元函数 $\Assign{f}{M}
  \colon \Domain{M}^n \to \Domain{M}$。
\end{enumerate}
\end{defn}

\begin{ex}
算术语言的一个结构~$\Struct M$ 包含一个集合（作为论域），一个元素 $\Assign{\Obj 0}{M}$ 作为常项符号~$\Obj 0$ 的解释，一个一元函数 $\Assign{\Obj \prime}{M} \colon \Domain{M} \to \Domain M$，两个二元函数 $\Assign{\Obj +}{M}$ 与 $\Assign{\Obj
    \times}{M}$（均为 $\Domain M^2 \to \Domain M$），以及一个二元关系 $\Assign{\Obj <}{M} \subseteq \Domain{M}^2$。

此类结构的一个显而易见的例子如下：
\begin{enumerate}
\item $\Domain N = \Nat$
\item $\Assign{\Obj 0}{N} = 0$
\item $\Assign{\Obj \prime}{N}(n) = n + 1$，对所有 $n \in \Nat$
\item $\Assign{\Obj +}{N}(n, m) = n + m$，对所有 $n, m \in \Nat$
\item $\Assign{\Obj \times}{N}(n, m) = n\cdot m$，对所有 $n, m \in \Nat$
\item $\Assign{\Obj <}{N} = \Setabs{\tuple{n, m}}{n \in \Nat, m \in
  \Nat, n < m}$
\end{enumerate}
如此定义的 $\Lang L_A$ 的结构~$\Struct N$ 被称为\emph{算术的标准模型}，因为它对~$\Lang L_A$ 中非逻辑常项的解释正如你所预期。

然而，对于~$\Lang L_A$ 还存在许多其他可能的结构。
例如，我们可以取整数集~$\Int$ 而非~$\Nat$ 作为论域，并相应地定义 $\Obj 0$、$\Obj \prime$、$\Obj +$、$\Obj \times$、$\Obj <$ 的解释。
但我们也可以定义与数字甚至毫无关联的~$\Lang L_A$ 的结构。
\end{ex}

\begin{ex}
集合论语言~$\Lang L_Z$ 的一个结构~$\Struct M$ 只需要一个集合和一个二元关系。
因此从技术上讲，例如，所有人的集合以及关系“$x$ 比 $y$ 年长”，或者 $\Nat$ 连同关系 $n \ge m$（$n, m \in \Nat$），都可以用作 $\Lang L_Z$ 的结构。

$\Lang L_Z$ 的一个特别有趣的结构是\emph{遗传有穷集}结构~$\Struct{{HF}}$：其论域中的元素本身就是集合，且 $\Obj \in$ 的解释确实就是关系“$x$ 是~$y$ 的元素”：
\begin{enumerate}
\item $\Domain{{{HF}}} = \emptyset \cup \Pow{\emptyset} \cup
  \Pow{\Pow{\emptyset}} \cup \Pow{\Pow{\Pow{\emptyset}}} \cup \dots$；
\item $\Assign{\Obj \in}{{{HF}}} = \Setabs{\tuple{x, y}}{x, y \in
  \Domain{{{HF}}}, x \in y}$。
\end{enumerate}
\end{ex}

\begin{digress}
我们对什么算作一个结构所作的规定，会影响我们的逻辑。
例如，禁止空论域的规定保证了（在关于量化语句的满足（或真）的通常定义下）$\lexists[x][(!A(x) \lor \lnot !A(x))]$ 是有效的——即，它是一个逻辑真理。
而规定所有常项符号必须指称论域中的一个对象，则保证了存在概括是一种可靠的推理模式：$!A(a)$，因此 $\lexists[x][!A(x)]$。
如果我们允许名字指称论域之外的对象，或不具有指称，那么我们就是在走向一种\emph{自由逻辑}；在自由逻辑中，存在概括需要一个额外的前提：$!A(a)$ 且 $\lexists[x][\eq[x][a]]$，因此 $\lexists[x][!A(x)]$。
\end{digress}

\end{document}
